\section{Классические коды коррекции ошибок (коды с повторением, линейные коды). Код Хэмминга.}

\subsection*{Коды с повторением}

Это простейший вид помехоустойчивого кодирования. 
\paragraph{Принцип:} Каждый бит сообщения $m$ передается $n$ раз.
\textit{Пример (n=3):} $0 \to 000, 1 \to 111$.
\begin{itemize}
    \item \textbf{Расстояние:} Минимальное кодовое расстояние $d_{min} = n$.
    \item \textbf{Исправляющая способность:} Позволяет исправлять $k = \lfloor \frac{n-1}{2} \rfloor$ ошибок методом «голосования» (мажоритарный принцип). 
    \item \textbf{Недостаток:} Крайне низкая скорость передачи данных (высокая избыточность). При $n=3$ скорость $R = 1/3$.
\end{itemize}

\subsection*{Линейные (групповые) коды}

Линейный код — это более эффективный способ внесения избыточности, основанный на линейной алгебре.

\paragraph{Определение}
Код называется \textbf{линейным}, если множество его кодовых слов образует подпространство в векторном пространстве $V_n$ над полем $\mathbb{Z}_2$.
\begin{itemize}
    \item \textbf{Свойство:} Сумма (по модулю 2) двух любых кодовых слов также является кодовым словом.
    \item \textbf{Нулевой вектор:} Нулевое слово всегда принадлежит линейному коду.
    \item \textbf{Расстояние:} Минимальное расстояние $d_{min}$ линейного кода равно минимальному весу (количеству единиц) ненулевого кодового слова.
\end{itemize}

\paragraph{Задание кода}
Линейный $(n, k)$-код (где $n$ — длина слова, $k$ — число инфо-бит) задается:
1. \textbf{Порождающей матрицей $G$:} размером $k \times n$. Кодовое слово $v$ получается как $v = m \cdot G$.
2. \textbf{Проверочной матрицей $H$:} размером $(n-k) \times n$. Для любого разрешенного слова $v$ выполняется $v \cdot H^T = \mathbf{0}$.

\subsection*{Код Хэмминга для исправления одного замещения}

Коды Хэмминга — это класс совершенных линейных кодов, которые исправляют ровно одну ошибку.

\paragraph{Принцип построения (на примере кода (7, 4))}
Мы добавляем к 4 информационным битам ($k=4$) 3 контрольных бита ($r=3$), итого длина $n=7$.
1. Контрольные биты ставятся на позиции, являющиеся степенями двойки (1, 2, 4).
2. Каждый контрольный бит отвечает за проверку на четность определенной группы позиций.
3. Группы выбираются так, чтобы в двоичной записи номера любой позиции «1» в $j$-м разряде означала, что эта позиция проверяется $j$-м контрольным битом.

\paragraph{Исправление ошибки (Синдром)}
При получении слова $v'$ вычисляется \textbf{синдром} $S = v' \cdot H^T$.
\begin{itemize}
    \item Если $S = \mathbf{0}$, ошибок нет.
    \item Если $S \neq \mathbf{0}$, то значение синдрома в двоичной системе счисления прямо указывает на \textbf{номер позиции} бита, в котором произошла ошибка (замещение).
\end{itemize}

\paragraph{Пример}
Если пришел вектор и синдром получился $110_2 = 6$, значит, нужно инвертировать 6-й бит в принятом сообщении, чтобы восстановить оригинал.

\paragraph{Преимущество:} Высокая скорость передачи (мало избыточных бит) при гарантированном исправлении одиночных ошибок.